\subsubsection{Subharmonic Functions}\label{subsec:subharmonic-functions}

We shall need a notion of subsolutions to the Dirichlet problem which does not require them to be of class $C^2(\Omega)$. The definition is motivated by the maximum principle.

\begin{Definition}
\label{def:continuous-subharmonic-function}
A function $u$ in $C^0(\overline\Omega)$ is called \textbf{subharmonic} (\textbf{superharmonic}), if for every ball $B$ with $\overline B\subset\Omega$ and every function $h\in C(\overline B)$ with $h$ harmonic in $B$ and $u\leq h$ ($u\geq h$) on $\partial B$, we have $u\leq h$ ($u\geq h$) in $B$. A subsolution (supersolution) of the Dirichlet problem is a function $u\in C(\overline\Omega)$ which is subharmonic (superharmonic) and such that $u\leq g$ ($u\geq g$) on $\partial\Omega$.
\end{Definition}

Clearly, if $\Delta u\geq0$, then $u$ is also subharmonic in the sense of the new definition. We note the following properties:

\begin{enumerate}
\item The strong maximum principle holds, i.e., if $u$ is subharmonic and $v$ is superharmonic with $v\geq u$ on $\partial\Omega$, then either $v>u$ in $\Omega$ or $v=u$ everywhere. We prove this by contradiction. Assume that $u-v$ assumes its maximum $M$ at some point $\bm{x}_0\in\Omega$, where $M\geq0$. If $u-v\equiv M$ throughout $\Omega$, it follows that $u=v$; hence we may assume that there are points in $\Omega$ where $u-v\not\equiv M$. In that case, we can choose $\bm{x}_0$ in such a way that there is a ball $B\subset\Omega$ centered at $\bm{x}_0$ such that $u-v$ does not equal $M$ on all of $\partial B$. Let $\bar u$ and $\bar v$ denote the harmonic functions on $B$ which are equal to $u$ and $v$, respectively, on $\partial B$. We find
\begin{equation}\label{eq:subharmonic-difference-at-maximum}
M=(u-v)(\bm{x}_0)\leq(\bar u-\bar v)(\bm{x}_0),
\end{equation}
and the right-hand side is strictly less than $M$ by the strong maximum principle for harmonic functions. Hence we have a contradiction. An immediate consequence is that every subsolution for the Dirichlet problem is less than or equal to every supersolution.\item Let $u$ be subharmonic in $\Omega$ and let $B$ be a ball with $\overline B\subset\Omega$. Let $\bar u$ be the harmonic function on $B$ satisfying $\bar u=u$ on $\partial B$. Then the function
\begin{equation}\label{eq:subharmonic-harmonic-lifting}
U(\bm{x})=\begin{cases}\bar u(\bm{x}),&x\in B\\u(\bm{x}),&x\in\Omega\setminus B\end{cases}
\end{equation}
is also subharmonic in $\Omega$ (cf. Problem~\ref{ex:subharmonic-harmonic-lifting}). $U$ is called the harmonic lifting of $u$ with respect to $B$.\item If $u_1,u_2,\ldots,u_N$ are subharmonic, then $\max\{u_1,u_2,\ldots,u_N\}$ is also subharmonic.\end{enumerate}
